\documentclass[11pt]{article}
\usepackage[margin=1in]{geometry}
\usepackage{amsmath,amssymb,amsthm,booktabs,microtype,hyperref}
\hypersetup{hidelinks,pdfauthor={T},pdftitle={Consent Guarantees for Multi-Step Assistance}}
\newtheorem{proposition}{Proposition}
\title{Consent Guarantees for Multi-Step Assistance}
\author{T}
\date{}
\begin{document}
\maketitle

\begin{abstract}
A model can learn to anticipate needs and still overstep when permission
is uncertain. We derive the maximal action set that satisfies a veto in
every situation compatible with an observation. A finite audit checks
28,672 combinations of observations and permissions. For multi-step
assistance, we characterize which permitted first moves retain a
guaranteed route to completion or safe handoff. A legal preparation
can fail this test when it creates a commitment that cannot be
canceled. A separate synthetic training comparison distinguishes
verified permission from imputed approval. An exact maximal-set
representation can shrink the continuation check when several
possible histories share the same available actions.
\end{abstract}

\section*{Summary and scope}
\paragraph{Question.} What may an assistant do when authority is
partly observed and later steps could require separate consent?
\paragraph{Finding.} The intersection of compatible permission
sets governs each effect. A finite-horizon continuation test
identifies which permitted first actions can still finish or hand
off safely in every compatible state.
\paragraph{Evidence and scope.} Exact finite checks cover
immediate permissions and multi-step continuations. The fitted
comparison uses synthetic supervision. The guarantees require
valid permission sets and transitions that include every feasible
authority change.

\section{Preference and authority}
Cooperative inverse reinforcement learning studies how a robot can
infer a human's objective from interaction \cite{cirl}. Inferring
that someone probably wants an outcome leaves a separate question
about who may choose how to pursue it. A development example should
make that distinction explicit. Compare offering to reschedule a
meeting with actually moving it. The latter also affects colleagues
and may require authorization the current user cannot give.

Let $C(o)$ be the nonempty set of states still compatible with a
trusted observation $o$. In state $s$, let $P(s)$ contain the actions
permitted by all relevant authorities under the declared policy.
Include separate actions for asking, deferring, and reversible
assistance. The set $P(s)$ can vary with jurisdiction and action scope.
\begin{proposition}
The maximal set of actions safe against any compatible veto is
\[
 P_*(o)=\bigcap_{s\in C(o)}P(s).
\]
A randomized action policy guarantees
this property exactly when its support lies in $P_*(o)$.
\end{proposition}
\begin{proof}
Each member of the intersection lies in every permission set. Each
excluded action has at least one compatible state in which it is
forbidden. Assigning that action positive probability gives positive
probability of violating a veto in that state. Conversely all mass
inside the intersection is permitted in each state.
\end{proof}

Place estimated benefit inside this permission boundary. Even
with correct gains, a positive weighted sum can compensate for
a protected person's loss by assigning sufficiently large gains
to others. The companion finite check enumerates 125 three-person
gain vectors. Of these, 27 have a positive sum and at least one
negative component. A protected welfare floor would reject those
vectors. Actual authorization is a separate input to $P(s)$,
and a predicted gain cannot provide it.

\section{Permission across a sequence of effects}
Suppose permitted preparatory actions can change the available
later choices. Let $T(s,a)$ be the nonempty set of possible
successor states after executing action $a$ in state $s$.
Let the resulting trusted observation $z$ reveal a subset of
those successors. Include authorization changes and any effects
bundled with the action in this model. A set $G$ consists of
states in which the workflow may finish or hand off safely.
Stopping in a state outside $G$ fails the stated completion
requirement.

For a current compatible set $B$, define $V_0(B)$ to hold
exactly when $B\subseteq G$. An action $a$ is viable with $h$
steps remaining if $h>0$, it is permitted throughout $B$, and
for each observation $z$ with a nonempty successor set
\[
B_{a,z}=\{s'\mid s'\in T(s,a)\text{ for some }s\in B,\ 
                       \text{the observation is }z\}
\]
we have $V_{h-1}(B_{a,z})$. Set $V_h(B)$ to hold when
$B\subseteq G$ or there is at least one viable action.
This is a finite-horizon use of the state-set construction
for games with imperfect information \cite{chatterjee}.

\begin{proposition}
The viable actions are exactly the first actions of policies
that obey all permissions and finish or hand off within $h$
steps in every modeled history. Under a pathwise guarantee,
a randomized policy can use only such first actions.
\end{proposition}
\begin{proof}
At horizon zero the only decision is to stop, which is safe
exactly when all compatible states are terminal. At positive
horizon the first action must satisfy the immediate permission
intersection. Every reachable observation leaves a continuation
problem with one fewer step. A legal policy needs a continuation
in each branch. Conversely, assembling one legal continuation
per branch after a permitted action gives a legal policy.
Induction proves the equivalence. A positive probability of a
nonviable action gives positive probability to a choice without
the required pathwise guarantee.
\end{proof}

Consider a document assistant that may either ask for consent
or prepare a release in an external tool. Asking reveals whether
it may commit, and a refusal permits safe handoff. Preparing is
currently permitted in both worlds but creates a reservation.
If consent is absent and the reservation has no permitted
cancel path, preparation can strand the workflow. With a
two-step horizon, only asking is viable. Adding a permitted
cancel action makes preparation viable too. The cancel path
must be enforced by the external tool, with a fresh permission
check for every later effect.

For 128 four-state fixtures, an independent oracle enumerated
all bounded policy trees and evaluated each in every compatible
initial world and successor. It agreed with the recurrence on
512 model and horizon combinations. Across the positive
horizons there were 341 immediately permitted first-action
occurrences with no viable continuation. The counts describe
the synthetic generator and do not estimate actual overreach.

\section{Checking viability with maximal compatible sets}
Viability is downward closed. A policy that works from $B$ also
works from any $B'\subseteq B$. Therefore it is enough to retain the
inclusion-maximal viable sets $\mathcal M_h$. The initial family is
$\mathcal M_0=\{G\}$. For a given action $a$, let $S_a$ be its set
of permitted states. For observation $z$ and a viable set $W$ at
horizon $h-1$, define
\[
 D_{a,z}(W)=\{s:\ T(s,a)\cap O_z\subseteq W\},
\]
where $O_z$ contains the states compatible with observation $z$.
For the deterministic observation map used in the checker, these
sets partition the state space.

\begin{proposition}
The inclusion-maximal beliefs permitting first action $a$ are the
maximal members of the family
\[
 \left\{S_a\cap\bigcap_z D_{a,z}(W_z):
             W_z\in\mathcal M_{h-1}\text{ for every }z\right\}.
\]
To obtain $\mathcal M_h$, take the maximal members of the union of
these action-specific families and $\{G\}$. The intersections can
be computed one observation at a time, removing contained sets
after each step.
\end{proposition}
\begin{proof}
An observation successor is viable exactly when it is contained
in some maximal viable set at the preceding horizon. That holds
for the successor after $a$ and $z$ precisely when every initial
state lies in the corresponding $D_{a,z}(W_z)$. Intersect with
$S_a$ for immediate permission. An observation with no reachable
successor imposes no additional constraint. This proves the family
formula. If $P\subseteq Q$ are two partial intersections, then
$P\cap D\subseteq Q\cap D$ for any later factor $D$. Dropping
$P$ therefore preserves the downward closure of the result.
Stopping contributes $G$ separately.
\end{proof}

Keep the maximal sets separately for each first action if the gate
must return every viable choice. A system can cache these sets for a
fixed permission model and look up a new belief by containment. A
changed permission or transition rule requires rebuilding the relevant
frontiers. This is an exact representation, with no general bound on
the number of maximal sets that a model can produce.

A structured fixture used sixteen possible states, eight probe
observations and four commit choices. Every state admitted the probe,
which revealed its pair. The correct commit then depended on that
pair. At the second horizon, direct expansion would enumerate
327,680 combinations across actions. The folded calculation used
160 pairwise intersections. This compares enumeration work on that
fixture. The check also compared all first actions against a subset
oracle across 131,072 belief and horizon pairs on the fixture and
9,216 on smaller models. One negative control sent two successor
states to the same observation. Each admitted a different terminal
move, and the gate rejected the shared predecessor because the
operator could not know which move was permitted.

\section{Labeling cases during training}
Construct matched triples with the same helpful objective: explicit
authorization, explicit veto, and missing or ambiguous authority.
Label the minimal legitimate act in each case. Preserve examples
where scope differs, such as permission to draft a message without
permission to send it. Train a permission prediction head with
evidence about who can authorize which effects. Use a structured
loss on probability mass assigned to excluded actions and a separate
loss rewarding practical, authorized help. A low average loss leaves
individual violations possible, so action execution still needs a
trusted, fresh authority check.

A permission query is valuable when it can distinguish the compatible
states and a permitted route remains open. Some requests are intrusive
or repetitive. Include cases where a reversible offer is a better
choice than asking again. A refusal carries no negative engagement
penalty in this curriculum.

For a multi-step curriculum, label each viable action from a
declared transition model. Keep tests where asking reveals a
refusal, authority expires between steps, or a previously offered
cancel path fails. If a transition is possible but missing from
the model, the guarantee is invalid. Treat approval evidence and
its freshness as inputs to the effect gate.

\section{Finite check and limits}
Three latent states, four actions, and every nonempty compatible-state
set give $16^3(2^3-1)=28{,}672$ checks. Bitwise intersection agreed
with a separately computed per-action veto scan. In the minimal
witness, an action permitted by one compatible state and vetoed by
another is excluded even if the permitted state is assigned high
probability. The guarantee requires the true state to be represented
and the permission sets to be correct. A fabricated authorization
or omitted affected party defeats those premises. A later study
should measure overreach and needless refusal separately, including
revocation and cross-person authority.
\section{Learning from imputed and verified permission}
Each seed trained two logistic policies with explicit approval,
veto, and unknown encoded separately. One training target treated
85\% of unknown cases as approved. The other required observed
affirmative authority. An independently sampled 3,200-case panel
contained 788 authorized cases and 2,412 cases with a veto or
unknown authority. The imputed target led to 1,641 unauthorized
executions. The verified target led to zero. Both executed in all
788 authorized cases. This binary experiment tests the supervision
rule for execution. It leaves the choice among asking, waiting, and
giving safe assistance to a larger action study.

\paragraph{Review status.} Independent manuscript review and a language-model study remain open.

\begin{thebibliography}{9}
\bibitem{cirl} D. Hadfield-Menell, A. Dragan, P. Abbeel and S. Russell.
Cooperative Inverse Reinforcement Learning.
\url{https://arxiv.org/abs/1606.03137}.
\bibitem{chatterjee} Krishnendu Chatterjee, Laurent Doyen,
Thomas A. Henzinger, and Jean-Fran\c{c}ois Raskin.
\emph{Algorithms for Omega-Regular Games with Imperfect Information}.
Logical Methods in Computer Science 3(3:4), 2007.
\url{https://arxiv.org/pdf/0706.2619}.
\end{thebibliography}
\end{document}
